% Why a constant release rate fails --- plan \S4.9 (Phase 3).
% Rewritten from the original Ch.4 \S9 attempt (preserved in the git-free
% folder as 09_potential_application_in_bmvr.tex). The old section proposed
% p = QS mean; this section states precisely why that cannot work, and ends
% at the door of the renewal formulation (\S\ref{sec:renewal}).

\section{From one cell to a population: why a constant release rate fails}
\label{sec:why-constant-fails}

The single-cell theory is complete. We now ask what it implies at the level
of a population of infected cells --- and begin with the standard object it
is meant to improve upon.

\subsection{The classical BMVR as comparator}
\label{sec:bmvr-classical}

In the simplified (target-cell-constant) Basic Model of Viral Replication
\cite{mclean1993balance,nowak1996population,perelson1996hiv}, infected
cells $\Icell$ and free particles $\Vfree$ obey
\begin{equation}
\ddt{\Icell} = \gamma T\,\Vfree - d_{\Icell}\,\Icell,
\qquad
\ddt{\Vfree} = p\,\Icell - c\,\Vfree,
\label{eq:BMVR-classical}
\end{equation}
with target cells $T$ held fixed, infection rate $\gamma$, virion clearance
$c$, and where the two constants that concern us are $p$, the per-cell
release rate, and $d_{\Icell}$, the infected-cell removal rate. Both are
phenomenological: they are parameters to be fitted, not quantities derived
from anything happening inside the cell. The basic reproduction number of
\eqref{eq:BMVR-classical} is
\begin{equation}
R_0=\frac{\gamma T\,p}{c\,d_{\Icell}}.
\end{equation}
Every cell, regardless of its age or its contents, releases at the same rate
$p$ and dies at the same rate $d_{\Icell}$. For a budding virus this is at
least a plausible picture. For a lytic pathogen it is plainly false: nothing
is released until the cell ruptures, and then everything is released at
once. The question of this part of the chapter is what
\eqref{eq:BMVR-classical} should be replaced with.

\subsection{The synchronous cohort}
\label{sec:sync-cohort}

Consider first a clean special case: $N_0$ cells, all infected at the same
instant, each running an independent copy of the single-cell process, with
no reabsorption or extracellular replication. The number still productively
infected at time $t$ is $N_0\,\Ihat(t)$ (exactly in expectation, and with
vanishing relative error as $N_0$ grows), and the expected cumulative
extracellular count is simply
\begin{equation}
\Vfree(t)=N_0\,V(t),
\end{equation}
since $V(t)$ is the expected release from one cell. In this synchronous
picture, the release \emph{rate} is $N_0 V'(t)=N_0\delta K(t)$: it depends
on the age of the cohort, and it is not constant. The classical model's
$p\Icell$ can mimic such a cohort only by accident, and only briefly.

\subsection{The proposal that cannot work}
\label{sec:proposal-fails}

The natural thought --- and the one originally pursued --- is to replace the
constant $p$ in \eqref{eq:BMVR-classical} by a constant built from the
intracellular theory: the quasi-stationary mean,
\begin{equation}
p\ \stackrel{?}{=}\ \langle X\rangle_{\rm QS}=\frac{a}{a-1}.
\label{eq:naive-p}
\end{equation}
The intuition is that a long-lived infected cell hovers at the QS level, so
its bursting should look like release at rate $\langle X\rangle_{\rm QS}$.
The obstruction is dimensional: $p$ has units of particles per unit time,
while $\langle X\rangle_{\rm QS}$ is a pure count. Substituting
\eqref{eq:naive-p} introduces a spurious factor of time into the model.
Worse, when $\mu=0$ one has $\langle X\rangle_{\rm QS}=V_\infty$ exactly, so
the proposal substitutes the \emph{total lifetime output} of a cell in place
of its \emph{output rate}.

Two repairs suggest themselves, and both fail for related reasons. First,
average the burst size over the burst-time density $\varphi=\delta J$ of
\S\ref{sec:burst-time}:
\begin{equation}
\int_0^\infty\varphi(t)\,\mean{\Kb\mid\tau=t}\,\dd t
=\delta\int_0^\infty K(t)\,\dd t
=V_\infty,
\end{equation}
using \eqref{eq:sizebias} and $V'=\delta K$ --- again a count, not a rate.
Second, weight the load-dependent flux by the burst-time density,
$\delta\int\varphi(t)\,\mean{\xt^2\mid\cdot}\,\dd t$: this has the right
dimensions, but it averages over the wrong measure. The burst-time density
describes \emph{when a single cell bursts}; what a population release rate
needs is an average over the \emph{ages of the cells currently infected}.
In an asynchronous infection those two distributions are not the same, and
no reweighting of a single-cell count by a single-cell density can fix the
substitution, because the difficulty is not in the single-cell theory at
all. It is in the population bookkeeping.

\subsection{What the attempt got right}
\label{sec:attempt-right}

Strip the failed proposal back and something correct remains. Suppose a
single infected macrophage founds a chain of infections: when it ruptures,
each released particle goes on to found a new intracellular process in a new
cell; these in turn end at different times, and the cycles continue,
beginning and ending, beginning and ending. After a sufficient time there
emerges a kind of stationarity: any given release event exhibits a burst
time and a burst size distributed according to the laws of
\S\ref{sec:burst-time-size}, and short-duration processes occur more
frequently than long ones simply because they turn over faster. What we want
is the overall rate of release events and their sizes --- the quantity that
should replace the $p\Icell$ term of \eqref{eq:BMVR-classical}.

Two observations from \S\ref{sec:sync-cohort}--\S\ref{sec:proposal-fails}
organise what follows. First, the expected release flux of a cell of age
$a$ is not a constant but $\delta K(a)$ --- the burst rate $\delta\xt$ times
the size $\xt$, averaged. Second, at any host time the infected population
is a mixture of ages, one cohort for each past moment of infection. The
correct population model must therefore keep track of \emph{cell age}: how
many cells were infected when, how many of them are still productive, and
how much each age class is releasing. That is precisely an age-structured,
or renewal, formulation, and because both of its ingredients already exist
in closed form, it is exactly solvable. We take it up next.
